\documentclass[11pt]{article}

\usepackage[margin=1in]{geometry}
\usepackage[T1]{fontenc}
\usepackage{lmodern}
\usepackage{microtype}
\usepackage{amsmath,amssymb,bm}
\usepackage{booktabs}
\usepackage{xcolor}
\usepackage{float}
\usepackage{hyperref}
\usepackage{enumitem}

\definecolor{accent}{RGB}{25,75,115}

\hypersetup{
  colorlinks=true,
  linkcolor=accent,
  citecolor=accent,
  urlcolor=accent
}

\newcommand{\dd}{\mathrm{d}}
\newcommand{\as}{a_s}
\newcommand{\ee}{\mathrm{e}}
\newcommand{\Jvec}{\bm J}
\newcommand{\Pmat}{\bm P}
\newcommand{\Mmat}{\bm M}

\title{\textbf{Validation of the Numerical \(b\)-Space Jet Evolution Solver}}
\author{}
\date{\today}

\begin{document}

\maketitle

\begin{abstract}
This report validates a numerical solver for the coupled quark--gluon
renormalization-group evolution of jet functions in impact-parameter space.
The running coupling and flavor thresholds are checked against analytic
relations, the time-like splitting kernels are compared with published Mellin
moments, and the discrete convolution is compared with direct continuum
quadrature. Exact manufactured solutions test forward and backward evolution
over the full numerical plane. An end-to-end calculation with physical
splitting kernels then measures grid and large-\(b\) closure effects.

The kernel moments agree with their published references to a maximum relative
difference of \(1.15\times10^{-4}\) through next-to-next-to-leading order
(NNLO). Euler and the second-order midpoint Runge--Kutta method (RK2) show
first- and second-order time convergence, respectively. The main numerical
limitation is the spatial convolution: with \(n_b=257\) spatial nodes, its relative difference
from continuum quadrature is \(2.89\%\)--\(6.17\%\), with approximately
first-order convergence. Full physical evolution is also sensitive to the
finite large-\(b\) closure, particularly in the backward-evolved NNLO gluon
channel.
\end{abstract}

\tableofcontents

\section{Evolution problem and numerical domain}

\subsection{Coupled \texorpdfstring{\(b\)}{b}-space evolution}

The solver evolves the quark-singlet and gluon jet functions in the
impact-parameter representation of the collinear energy--energy correlator
(EEC) jet factor. The variable \(b\), measured in \(\mathrm{GeV}^{-1}\), is
Fourier conjugate to the transverse momentum entering the measurement. The
two functions are arranged as the row vector
\[
  \Jvec(b,\mu)=\bigl(J_q(b,\mu),J_g(b,\mu)\bigr),
\]
according to
\begin{equation}
  \frac{\partial}{\partial\ln\mu^2}\Jvec(b,\mu)
  =
  \int_0^1 \dd y\,y^2\,
  \Jvec\left(\frac{b}{y},\mu\right)\Pmat(y,\mu),
  \label{eq:continuum-rg}
\end{equation}
where
\[
  \Pmat=
  \begin{pmatrix}
    P_{qq} & P_{qg}\\
    P_{gq} & P_{gg}
  \end{pmatrix}.
\]
The convention is that \(P_{ij}\) describes a daughter of type \(i\)
produced from a parent of type \(j\). The component equations are therefore
\begin{align}
  \partial_tJ_q
  &=
  \int_0^1\dd y\,y^2
  \left[
    J_q(b/y,t)P_{qq}(y,t)
    +J_g(b/y,t)P_{gq}(y,t)
  \right],
  \label{eq:component-q}\\
  \partial_tJ_g
  &=
  \int_0^1\dd y\,y^2
  \left[
    J_q(b/y,t)P_{qg}(y,t)
    +J_g(b/y,t)P_{gg}(y,t)
  \right],
  \label{eq:component-g}
\end{align}
with \(t=\ln\mu^2\). The flavor multiplicity of the quark-singlet channel is
included in \(P_{qg}\).

The perturbative expansion is
\begin{equation}
  P_{ij}(y,\mu)
  =
  \sum_{L=0}^{\infty}
  \as^{L+1}(\mu)P_{ij}^{(L)}(y),
  \qquad
  \as=\frac{\alpha_s}{4\pi}.
  \label{eq:pert-series}
\end{equation}
The solver options \texttt{order=0}, \(1\), and \(2\) retain the cumulative
leading-order (LO), next-to-leading-order (NLO), and NNLO series,
respectively.

\subsection{Boundary curve and lattice}

Boundary data are supplied on the canonical scale
\begin{equation}
  \mu_*(b)=\frac{b_0}{b_*(b)},
  \qquad
  t_*(b)=\ln\mu_*^2(b).
  \label{eq:boundary-curve}
\end{equation}
Here \(b_0=2\ee^{-\gamma_E}\simeq1.12292\) is dimensionless, while \(b_*(b)\)
has the same units as \(b\). A standard saturating prescription is
\[
  b_*(b)=\frac{b}{\sqrt{1+b^2/b_{\mathrm{sat}}^2}},
\]
which approaches \(b\) for \(b\ll b_{\mathrm{sat}}\) and
\(b_{\mathrm{sat}}\) for \(b\gg b_{\mathrm{sat}}\). The solver accepts any
positive prescription whose sampled canonical scales have the ordering
described below.

Choose \(n_b\ge2\) spatial nodes, uniform in \(\ell=\ln b\) and ordered from
large to small \(b\):
\begin{align}
  \ell_i&=\ell_{\max}-(i-1)\Delta\ell,
  &
  \Delta\ell&=\frac{\ell_{\max}-\ell_{\min}}{n_b-1},
  \\
  b_i&=\ee^{\ell_i},
  &
  t_i&=t_*(b_i),
  \qquad i=1,\ldots,n_b.
\end{align}
The \(t\)-grid is therefore derived entirely from the same \(b\)-grid and
\(b_*\) prescription: there is exactly one \(t_i\) for every \(\ell_i\), with
no independent time range or sorting. The implementation requires
\(t_{i+1}>t_i\). Boundary values occupy the diagonal lattice points
\((\ell_i,t_i)\) exactly. At fixed \(b_i\), entries with time index \(n>i\)
are obtained by forward evolution and those with \(n<i\) by backward
evolution.

The convolution variable is \(u=-\ln y\), sampled with the same spacing,
\(u_m=m\Delta\ell\) and \(y_m=\ee^{-u_m}\). The identity
\[
  \ln(b_i/y_m)=\ell_i+u_m=\ell_{i-m}
\]
places every shifted argument on another spatial node. This avoids
interpolation inside the convolution and gives the discrete dependence its
one-sided triangular structure.

The numerical domain is finite. Whenever \(b/y>b_{\max}\), the solver uses
\begin{equation}
  \Jvec(b>b_{\max},t)=0.
  \label{eq:zero-closure}
\end{equation}
This closure is appropriate only when the physical boundary and evolved
solution are already negligible near \(b_{\max}\). Otherwise the omitted
large-\(b\) tail becomes a source of numerical error, especially in diagonal
channels containing plus distributions.

\subsection{Dimensionless comparison measures}

Scalar quantities with a nonzero reference value are compared using
\begin{equation}
  \delta_{\mathrm{rel}}
  =
  \left|
    \frac{x_{\mathrm{num}}}{x_{\mathrm{ref}}}-1
  \right|.
  \label{eq:scalar-relative}
\end{equation}
For arrays, two relative measures are used:
\begin{align}
  \epsilon_2
  &=
  \frac{\lVert \bm x_{\mathrm{num}}-\bm x_{\mathrm{ref}}\rVert_2}
       {\lVert \bm x_{\mathrm{ref}}\rVert_2},
  \label{eq:relative-l2}\\
  \epsilon_{\mathrm{RMS}}
  &=
  \sqrt{
    \frac{1}{n_{\mathrm{samp}}}
    \sum_k
    \left|
      \frac{x_k^{\mathrm{num}}}{x_k^{\mathrm{ref}}}-1
    \right|^2
  }.
  \label{eq:pointwise-rms}
\end{align}
The first is a global relative \(L^2\) difference; the second gives equal
weight to the pointwise relative difference at each sampled value.

When two resolutions are compared, the observed convergence order is
\begin{equation}
  p=
  \frac{\ln(\epsilon_{\mathrm{coarse}}/\epsilon_{\mathrm{fine}})}
       {\ln(h_{\mathrm{coarse}}/h_{\mathrm{fine}})}.
  \label{eq:observed-order}
\end{equation}
Here \(h=\Delta\ell\) for the spatial-convolution test and \(h=\Delta t\), or
the maximum \(\Delta t\) for a nonuniform grid, for the time-evolution tests.

Some exact identities require two nonzero contributions to cancel. A
relative difference from zero is undefined, so cancellation is measured by
the normalized residual
\begin{equation}
  \eta(a,b)=\frac{|a+b|}{|a|+|b|}.
  \label{eq:normalized-residual}
\end{equation}
This is dimensionless and measures the defect relative to the terms being
cancelled. All result tables below use relative differences or normalized
residuals.

\section{Active-flavor prescription}

Two flavor schemes are supported: a variable-flavor-number scheme (VFNS) and
a fixed five-flavor scheme:
\begin{equation}
n_f(\mu;\texttt{:VFNS})=
\begin{cases}
3,&\mu<m_c,\\
4,&m_c\le\mu<m_b,\\
5,&m_b\le\mu<m_t,\\
6,&m_t\le\mu,
\end{cases}
\qquad
n_f(\mu;\texttt{:nf5})=5.
\label{eq:nf-schemes}
\end{equation}
The current threshold values are
\[
  m_c=1.27~\mathrm{GeV},\qquad
  m_b=4.18~\mathrm{GeV},\qquad
  m_t=172.76~\mathrm{GeV}.
\]

The tests evaluate the flavor function immediately below and exactly at every
threshold. They also construct splitting-kernel table rows on both sides of a
threshold and compare them with direct kernel evaluations using the flavor
number selected at the same scale. Both schemes return the specified values,
the kernel tables use the same prescription, and invalid schemes or
nonpositive scales are rejected.

These checks establish consistency of flavor selection within the numerical
implementation. The threshold masses are fixed input parameters for this
validation; their phenomenological determination is outside its scope.

\section{Numerical strong coupling}

Within an interval of fixed \(n_f\), the one-loop running coupling satisfies
\begin{equation}
  \frac{1}{\alpha_s(\mu_2)}
  =
  \frac{1}{\alpha_s(\mu_1)}
  +\frac{\beta_0(n_f)}{2\pi}
   \ln\frac{\mu_2}{\mu_1}.
  \label{eq:one-loop-alpha}
\end{equation}
The numerical coupling is initialized from
\[
  \alpha_s(M_Z)=0.118,\qquad M_Z=91.1876~\mathrm{GeV}.
\]
It integrates the \(K\)-loop beta function with the Tsitouras 5/4 Runge--Kutta
method, using relative and absolute ODE tolerances \(10^{-10}\) and
\(10^{-12}\). The analytic comparison below sets \(K=1\); the subsequent
fixed-flavor comparison sets \(K=4\).

For VFNS, this expression is applied successively across every flavor
interval between the reference and target scales. Numerical one-loop
evolution agrees with this piecewise analytic result to better than
\(2\times10^{-8}\) relative at
\[
  \mu=1.1,\ 2,\ 10,\ 300~\mathrm{GeV}.
\]
The tested coupling is positive, decreases monotonically with \(\mu\), and
is continuous across \(m_c\), \(m_b\), and \(m_t\).

As a separate fixed-\(n_f=5\) regression, the four-loop result was compared
with an independently integrated ODE using stored rounded beta coefficients
and with a closed-form truncated resummed expansion. All three calculations
use the common reference pair
\(\alpha_s(91.2~\mathrm{GeV})=0.118\) and are evaluated over
\[
  \mu=2,\ 4.18,\ 5,\ 10,\ 20,\ 50,\ 91.2,\ 100,\ 160,\ 500~\mathrm{GeV}.
\]

\begin{table}[H]
\centering
\caption{Maximum relative difference between the numerical
\texttt{:nf5} coupling and two independent fixed-flavor references.}
\label{tab:fixed-alpha}
\begin{tabular}{lr}
\toprule
Reference & Maximum relative difference\\
\midrule
Numerical ODE with stored rounded coefficients & \(2.22\times10^{-7}\)\\
Truncated resummed expression & \(1.59\times10^{-3}\)\\
\bottomrule
\end{tabular}
\end{table}

The first difference is consistent with replacing rounded stored beta
coefficients by analytic \(n_f\)-dependent coefficients. The larger
difference from the truncated expression occurs at \(2~\mathrm{GeV}\).
All formulations are equal by construction at the shared reference scale.

The VFNS implementation changes beta-function coefficients at the heavy-quark
thresholds while keeping \(\alpha_s\) continuous. It does not include finite
higher-order decoupling coefficients, so it represents continuous threshold
switching rather than a complete precision matching prescription.

\section{Splitting kernels and published Mellin moments}

For a splitting-kernel coefficient, define the Mellin moment
\begin{equation}
  \gamma_{T,ij}^{(L)}(N)
  =
  -\int_0^1\dd y\,y^{N-1}P_{ij}^{(L)}(y).
  \label{eq:published-moment}
\end{equation}
For a diagonal channel, this integral combines the regular function, delta
function, and plus-distribution terms. It therefore probes the complete
distributional kernel rather than isolated values of its regular part.

Appendix A of Ref.~\cite{Dixon:2019uzg} gives \(N=3\) time-like moments in
the convention of Eq.~\eqref{eq:published-moment}. All four singlet channels
were compared through NNLO for \(n_f=3,4,5,6\). The same reference also
provides moments with an additional \(\ln y\) weight and, at LO, a
\(\ln^2y\) weight. Because these logarithmic weights vanish at \(y=1\), they
test combinations of regular and plus-distribution terms that differ from
the ordinary \(N=3\) moments.

\begin{table}[H]
\centering
\caption{Maximum relative difference from the published moments over all
four channels and \(n_f=3,4,5,6\).}
\label{tab:published-moments}
\begin{tabular}{lr}
\toprule
Quantity & Maximum relative difference\\
\midrule
\(P^{(0)}\), LO & \(4.44\times10^{-16}\)\\
\(P^{(1)}\), NLO & \(4.23\times10^{-10}\)\\
\(P^{(2)}\), NNLO & \(1.15\times10^{-4}\)\\
Published logarithmic moments & \(2.29\times10^{-10}\)\\
\bottomrule
\end{tabular}
\end{table}

The LO result is at floating-point roundoff, and the NLO and logarithmic
moments agree at numerical integration precision. The largest NNLO
difference occurs in the \(qg\) moment at \(n_f=6\). Its larger size is
expected because the three-loop kernels are implemented as decimal
parametrizations rather than exact analytic functions.

\section{Conservation of longitudinal momentum}

In a time-like collinear splitting, \(y\) is the fraction of the parent's
large light-cone momentum carried by a daughter parton. For massless
collinear partons this is also the daughter energy fraction. Conservation of
this quantity requires the first momentum moment of all daughters from a
fixed parent to sum to zero.

Define
\begin{equation}
  M_{ij}
  =
  \int_0^1\dd y\,y\,P_{ij}(y).
  \label{eq:momentum-matrix}
\end{equation}
The quark-parent and gluon-parent identities are
\begin{align}
  M_{qq}+M_{gq}&=0,
  \label{eq:q-parent-sum}\\
  M_{qg}+M_{gg}&=0.
  \label{eq:g-parent-sum}
\end{align}
The quark multiplicity is already included in the singlet \(P_{qg}\).
Equations~\eqref{eq:q-parent-sum} and \eqref{eq:g-parent-sum} are the
\(N=2\) Mellin-moment relations because the Mellin weight is
\(y^{N-1}=y\).

Real-emission terms redistribute momentum between daughter channels, while
the delta and plus-distribution terms supply the virtual contribution needed
for the cancellation. The quality of that cancellation is measured with
Eq.~\eqref{eq:normalized-residual}:
\[
  \eta_q=\eta(M_{qq},M_{gq}),
  \qquad
  \eta_g=\eta(M_{qg},M_{gg}).
\]

\begin{table}[H]
\centering
\caption{Normalized momentum-conservation residuals for each perturbative
coefficient at \(n_f=5\).}
\label{tab:coefficient-momentum}
\begin{tabular}{lrr}
\toprule
Coefficient & Quark parent \(\eta_q\) & Gluon parent \(\eta_g\)\\
\midrule
\(P^{(0)}\) & \(6.25\times10^{-17}\) & \(0\)\\
\(P^{(1)}\) & \(2.10\times10^{-10}\) & \(2.35\times10^{-10}\)\\
\(P^{(2)}\) & \(1.64\times10^{-6}\) & \(4.61\times10^{-6}\)\\
\bottomrule
\end{tabular}
\end{table}

The solver uses the cumulative perturbative series rather than an isolated
coefficient. Table~\ref{tab:cumulative-momentum} gives the corresponding
normalized residuals at \(\as=0.01\).

\begin{table}[H]
\centering
\caption{Normalized momentum-conservation residuals of the cumulative
kernels at \(\as=0.01\) and \(n_f=5\).}
\label{tab:cumulative-momentum}
\begin{tabular}{lrr}
\toprule
Included terms & Quark parent \(\eta_q\) & Gluon parent \(\eta_g\)\\
\midrule
LO & \(0\) & \(0\)\\
LO+NLO & \(1.67\times10^{-11}\) & \(3.57\times10^{-11}\)\\
LO+NLO+NNLO & \(1.07\times10^{-8}\) & \(9.56\times10^{-8}\)\\
\bottomrule
\end{tabular}
\end{table}

The coefficient-level residual grows at NNLO because of the decimal
three-loop parametrization, but it remains below five parts per million.
Perturbative weighting reduces the residual of the cumulative NNLO kernel
below \(10^{-7}\) for \(\as=0.01\).

\section{Discrete physical convolution}

The discrete convolution must combine the regular, delta, and generalized
plus-distribution pieces with the correct channel orientation. It is tested
using the smooth functions
\begin{align}
  J_q(b)&=\ee^{-2b}(1+0.3b),\\
  J_g(b)&=0.7\,\ee^{-1.5b}(1+0.2b),
\end{align}
which are negligible before the closure point \(b_{\max}=30\). This prevents
the finite-domain boundary from dominating the comparison.

At the midpoint node,
\[
  i=(n_b+1)/2,\qquad
  b_i=\sqrt{b_{\min}b_{\max}}
      =0.173205\ldots~\mathrm{GeV}^{-1},
\]
one application of the lattice right-hand side is compared with adaptive
continuum quadrature. The continuum reference
evaluates each plus prescription directly as the difference between the
shifted real value and the endpoint value; it does not use the lattice
weights. The calculation uses
\[
  n_b=65,\ 129,\ 257,\qquad
  b_{\min}=10^{-3}~\mathrm{GeV}^{-1},\qquad
  b_{\max}=30~\mathrm{GeV}^{-1},\qquad
  \mu=20~\mathrm{GeV},
\]
the VFNS coupling evolved with the four-loop beta function, and cumulative
LO, NLO, and NNLO kernels.

\begin{table}[H]
\centering
\caption{Relative convolution difference at \(n_b=257\) and observed
convergence order between \(n_b=129\) and \(n_b=257\).}
\label{tab:physical-convolution}
\begin{tabular}{lrrrr}
\toprule
Perturbative order
  & Quark difference & Quark rate
  & Gluon difference & Gluon rate\\
\midrule
LO & \(6.002\%\) & 0.978 & \(4.860\%\) & 0.987\\
NLO & \(6.174\%\) & 0.964 & \(3.358\%\) & 1.292\\
NNLO & \(6.162\%\) & 0.964 & \(2.892\%\) & 1.421\\
\bottomrule
\end{tabular}
\end{table}

Both channels improve under each refinement at every perturbative order.
The channel mapping and distributional formula therefore converge toward
the independent continuum calculation. The observed rates also identify the
dominant limitation: the quark convolution is approximately first order and
  still differs from the continuum result by about \(6\%\) at \(n_b=257\).
The larger two-grid rates in the gluon channel can include cancellation
between separate error contributions and should not be interpreted as an
asymptotic order above one.

\section{Local evolution across the boundary curve}

Time integration is isolated from the nonlocal convolution by choosing
\begin{equation}
  \Pmat(y)=
  \begin{pmatrix}
    c_q & 0\\
    0 & c_g
  \end{pmatrix}
  \delta(1-y).
  \label{eq:local-kernel}
\end{equation}
Each spatial node then obeys an independent ordinary differential equation.
For boundary data at \(t_*(b)\), the exact solution on either side of the
boundary is
\begin{align}
  J_q(b,t)
  &=J_q\bigl(b,t_*(b)\bigr)
    \exp\left[c_q\bigl(t-t_*(b)\bigr)\right],\\
  J_g(b,t)
  &=J_g\bigl(b,t_*(b)\bigr)
    \exp\left[c_g\bigl(t-t_*(b)\bigr)\right].
\end{align}

The calculation uses \(0.5\le b/(\mathrm{GeV}^{-1})\le2.0\),
\(b_*(b)=b\),
\[
  J_q\bigl(b,t_*(b)\bigr)=1+0.2b,\qquad
  J_g\bigl(b,t_*(b)\bigr)=0.7+0.1b,
\]
and rates \(c_q=0.4\), \(c_g=-0.3\). Every stored value in both full-plane
regions is compared with the exact exponential.

\begin{table}[H]
\centering
\caption{Relative convergence of the local full-plane evolution.}
\label{tab:local-results}
\begin{tabular}{llrrrr}
\toprule
Method & \(n_b\) & \(\Delta t\)
  & Relative RMS & Maximum relative & Observed order\\
\midrule
Euler & 17  & \(1.733\times10^{-1}\) & \(1.340\times10^{-2}\) & \(3.951\times10^{-2}\) & --\\
Euler & 33  & \(8.664\times10^{-2}\) & \(6.531\times10^{-3}\) & \(1.948\times10^{-2}\) & 1.036\\
Euler & 65  & \(4.332\times10^{-2}\) & \(3.224\times10^{-3}\) & \(9.674\times10^{-3}\) & 1.018\\
Euler & 129 & \(2.166\times10^{-2}\) & \(1.602\times10^{-3}\) & \(4.821\times10^{-3}\) & 1.009\\
\midrule
RK2 & 17  & \(1.733\times10^{-1}\) & \(2.959\times10^{-4}\) & \(9.359\times10^{-4}\) & --\\
RK2 & 33  & \(8.664\times10^{-2}\) & \(7.175\times10^{-5}\) & \(2.279\times10^{-4}\) & 2.044\\
RK2 & 65  & \(4.332\times10^{-2}\) & \(1.766\times10^{-5}\) & \(5.623\times10^{-5}\) & 2.022\\
RK2 & 129 & \(2.166\times10^{-2}\) & \(4.382\times10^{-6}\) & \(1.397\times10^{-5}\) & 2.011\\
\bottomrule
\end{tabular}
\end{table}

Euler converges at first order and midpoint RK2 at second order. Because the
comparison covers both sides of the boundary, it verifies the time-step sign,
boundary indexing, and preservation of diagonal boundary data in both the
forward and backward passes.

\section{Nonlocal full-plane evolution}

The local calculation does not exercise shifted arguments or channel mixing.
A second manufactured problem collects the semi-discrete state into
\[
  \bm x(t)=
  \bigl(J_{q,1},\ldots,J_{q,n_b},
        J_{g,1},\ldots,J_{g,n_b}\bigr)^T
\]
and evolves it with
\begin{equation}
  \frac{\dd\bm x}{\dd t}
  =
  \left[1+\lambda(t-t_0)\right]\bm A_0\bm x,
  \qquad
  \lambda=0.15.
  \label{eq:manufactured-ode}
\end{equation}
The operator is specified component by component as
\begin{align}
  (\bm A_0\bm x)_{q,i}
  &=
  0.12J_{q,i}
  +\mathbf{1}_{i>1}
  \left(0.07J_{q,i-1}+0.09J_{g,i-1}\right),
  \\
  (\bm A_0\bm x)_{g,i}
  &=
  -0.08J_{g,i}
  +\mathbf{1}_{i>1}
  \left(-0.05J_{q,i-1}+0.04J_{g,i-1}\right).
  \label{eq:manufactured-operator}
\end{align}
It contains local and one-node-shift terms in all four channels, with the
same one-sided dependence as the lattice convolution. In the channel-major
ordering of \(\bm x\), \(\bm A_0\) is a \(2\times2\) block matrix whose four
spatial blocks are lower bidiagonal. At
\(t_0=t_1\), the reference state is
\begin{align}
  J_{q,i}(t_0)&=1+0.2\,b_i/b_{\max},\\
  J_{g,i}(t_0)&=0.8-0.1\,b_i/b_{\max}.
  \label{eq:manufactured-initial}
\end{align}

Because the time dependence is a scalar multiple of a fixed matrix, the exact
solution is
\begin{equation}
  \bm x(t)
  =
  \exp\left\{
    \left[
      \Delta t+\frac{\lambda}{2}\Delta t^2
    \right]\bm A_0
  \right\}\bm x(t_0).
  \label{eq:manufactured-exact}
\end{equation}
The test uses
\[
  0.2\le b/(\mathrm{GeV}^{-1})\le2.0,\qquad
  b_*(b)=\frac{b}{\sqrt{1+(b/b_{\mathrm{sat}})^2}},
  \qquad b_{\mathrm{sat}}=0.8~\mathrm{GeV}^{-1}.
\]
At every diagonal lattice point, the boundary input is taken directly from
Eq.~\eqref{eq:manufactured-exact}:
\[
  J_{q,i}(t_i)=[\bm x(t_i)]_{q,i},
  \qquad
  J_{g,i}(t_i)=[\bm x(t_i)]_{g,i}.
\]
The numerical method then reconstructs both full-plane regions without using
the exact solution away from this boundary.

\begin{table}[H]
\centering
\caption{Global relative \(L^2\) difference for the manufactured nonlocal
problem.}
\label{tab:nonlocal-results}
\begin{tabular}{lrrr}
\toprule
Method & \(n_b=17\) & \(n_b=33\) & Observed order\\
\midrule
Euler & \(1.0754\times10^{-2}\) & \(5.2908\times10^{-3}\) & 1.031\\
RK2 & \(1.9121\times10^{-4}\) & \(4.7539\times10^{-5}\) & 2.022\\
\bottomrule
\end{tabular}
\end{table}

The expected first- and second-order rates are recovered. This calculation
simultaneously exercises all four local and shifted channel couplings,
nonuniform time steps inherited from the curved boundary, midpoint kernel
evaluation, forward marching to \(n>i\), and reverse marching to \(n<i\).

\section{End-to-end evolution with physical kernels}

\subsection{Stationary continuum solution}

The momentum identity provides a simple exact solution of the continuum
equation. If
\[
  J_q(b,\mu)=J_g(b,\mu)=b,
\]
then \(J_i(b/y)=b/y\), and the \(y^2\) weight in
Eq.~\eqref{eq:continuum-rg} leaves the momentum weight \(y\):
\begin{equation}
  \partial_t\bigl[b(1,1)\bigr]
  =
  b(1,1)\int_0^1\dd y\,y\,\Pmat(y)=0.
  \label{eq:stationary-from-momentum}
\end{equation}
Thus
\begin{equation}
  J_q(b,\mu)=J_g(b,\mu)=b
  \label{eq:stationary-solution}
\end{equation}
is independent of \(\mu\) in the continuum.

This calculation uses the numerical running coupling, physical splitting
kernels, physical convolution, and full forward and backward evolution. The
supplied boundary condition is
\[
  J_q\bigl(b,t_*(b)\bigr)=J_g\bigl(b,t_*(b)\bigr)=b.
\]
The active-flavor prescription is VFNS and the coupling uses the four-loop
beta function.
However, the exact mode grows with \(b\), whereas the numerical closure sets
it abruptly to zero beyond \(b_{\max}\). The test is therefore an
end-to-end diagnostic of the implementation and its finite-domain
sensitivity, not a pure measurement of the time integrator.

\subsection{Baseline result}

The end-to-end tables show LO and NNLO, the two endpoints of the supported
perturbative truncations. NLO is already exercised independently in the
kernel-moment and physical-convolution tests; including it here would not add
a new numerical pathway. Results are summarized for
\(b\le0.03~\mathrm{GeV}^{-1}\), a window at least a factor of \(10^3\) below
the baseline closure point. This choice reduces direct endpoint
contamination while retaining points on both sides of the boundary curve; the
remaining closure sensitivity is measured separately below.

The baseline grid uses
\[
  b_{\min}=0.001~\mathrm{GeV}^{-1},\qquad
  b_{\max}=30~\mathrm{GeV}^{-1},\qquad
  n_b=1000,
\]
midpoint RK2, and
\[
  b_*(b)=\frac{b}{\sqrt{1+b^2/b_{\mathrm{sat}}^2}},
  \qquad
  b_{\mathrm{sat}}=1.12292~\mathrm{GeV}^{-1}.
\]
Table~\ref{tab:stationary-1000} summarizes points in the stated small-\(b\)
window.
The upper and lower regions lie above and below the boundary curve,
respectively.

\begin{table}[H]
\centering
\caption{Relative deviation from the stationary solution at \(n_b=1000\) and
\(b\le0.03~\mathrm{GeV}^{-1}\).}
\label{tab:stationary-1000}
\begin{tabular}{llrrrr}
\toprule
& & \multicolumn{2}{c}{Quark} & \multicolumn{2}{c}{Gluon}\\
\cmidrule(lr){3-4}\cmidrule(lr){5-6}
Order & Region & RMS & Maximum & RMS & Maximum\\
\midrule
LO & Upper & \(0.242\%\) & \(0.966\%\) & \(0.599\%\) & \(2.229\%\)\\
LO & Lower & \(0.395\%\) & \(0.613\%\) & \(1.114\%\) & \(2.120\%\)\\
NNLO & Upper & \(0.075\%\) & \(0.183\%\) & \(0.057\%\) & \(0.113\%\)\\
NNLO & Lower & \(0.495\%\) & \(1.336\%\) & \(3.036\%\) & \(12.335\%\)\\
\bottomrule
\end{tabular}
\end{table}

The largest deviation is in the backward-evolved NNLO gluon channel. The
separate \(n_f=5\), \(\as=0.01\) benchmark gives a normalized cumulative
kernel residual below \(10^{-7}\), but it does not bound that residual over
the running-coupling VFNS plane. The several-percent deviation is consistent
with the combined effects of spatial discretization, closure, time
resolution, repeated evolution, and kernel parametrization; these
contributions are not isolated by this test.

\subsection{Grid refinement}

\begin{table}[H]
\centering
\caption{Lower-region relative RMS deviation for
\(b\le0.03~\mathrm{GeV}^{-1}\) under grid refinement.}
\label{tab:lower-refinement}
\begin{tabular}{lrrrr}
\toprule
\(n_b\) & LO quark & LO gluon & NNLO quark & NNLO gluon\\
\midrule
500 & \(0.761\%\) & \(2.343\%\) & \(0.980\%\) & \(2.949\%\)\\
1000 & \(0.395\%\) & \(1.114\%\) & \(0.495\%\) & \(3.036\%\)\\
2000 & \(0.225\%\) & \(0.558\%\) & \(0.353\%\) & \(2.840\%\)\\
\bottomrule
\end{tabular}
\end{table}

The LO deviations decrease approximately with lattice spacing, and the NNLO
quark result also improves. The NNLO gluon result changes more slowly and is
not monotone between 500 and 1000 nodes. This behavior indicates that more
than one error source contributes in the backward-evolved NNLO region.

\subsection{\texorpdfstring{Large-\(b\)}{Large-b} closure sensitivity}

The closure point was moved from \(30\) to \(300~\mathrm{GeV}^{-1}\) while
approximately preserving \(\Delta\ell\):
\[
  (n_b,b_{\max}/\mathrm{GeV}^{-1})=(1000,30),\qquad(1223,300).
\]

\begin{table}[H]
\centering
\caption{Relative RMS deviation for \(b\le0.03~\mathrm{GeV}^{-1}\) after
moving the closure point at approximately fixed spatial resolution.}
\label{tab:closure-comparison}
\begin{tabular}{llrr}
\toprule
Region and order & Channel & \(b_{\max}=30\) & \(b_{\max}=300\)\\
\midrule
LO upper & Quark & \(0.242\%\) & \(0.039\%\)\\
LO upper & Gluon & \(0.599\%\) & \(0.050\%\)\\
NNLO lower & Quark & \(0.495\%\) & \(0.898\%\)\\
NNLO lower & Gluon & \(3.036\%\) & \(2.074\%\)\\
\bottomrule
\end{tabular}
\end{table}

Moving the closure ten times farther away strongly improves the LO upper
region and reduces the NNLO lower gluon deviation. It worsens the NNLO lower
quark result, showing that closure is important but is not the only source of
the backward-region discrepancy.

\section{Overall assessment}

The validation establishes the following:
\begin{itemize}[leftmargin=1.5em]
  \item The VFNS and fixed-\(n_f=5\) schemes select and propagate the intended
  active-flavor count.
  \item Numerical running agrees with piecewise analytic one-loop evolution
  and with an independent fixed-flavor ODE using rounded beta coefficients.
  \item All four splitting channels reproduce published ordinary and
  logarithmic Mellin moments through the accuracy of the NNLO
  parametrizations.
  \item At \(n_f=5\), longitudinal-momentum conservation holds with normalized
  residuals below \(5\times10^{-6}\) coefficient by coefficient and below
  \(10^{-7}\) for the cumulative NNLO kernel at \(\as=0.01\).
  \item The convolution converges toward independent continuum quadrature,
  but its approximately first-order spatial error is larger than the
  nominal RK2 time error at the tested resolutions.
  \item Forward and backward time evolution recover the expected Euler and
  RK2 convergence rates for both local and nonlocal exact problems.
  \item End-to-end physical evolution is most sensitive in the
  backward-evolved NNLO gluon channel and requires explicit checks of both
  grid resolution and the large-\(b\) endpoint.
\end{itemize}

The principal limitations are:
\begin{itemize}[leftmargin=1.5em]
  \item The present spatial convolution is approximately first order.
  \item A sharp zero closure is unsuitable if the jet function is not already
  negligible at \(b_{\max}\), particularly for diagonal plus distributions.
  \item Refining \(n_b\) changes spatial and time resolution simultaneously,
  which prevents a clean separation of their errors in the physical solver.
  \item VFNS uses continuous threshold switching without finite higher-order
  decoupling corrections.
  \item No independent exact solution is currently available for the full
  physical equation with a decaying boundary and a running coupling.
\end{itemize}

\section{Reproducibility}

The calculations were run with Julia 1.11.6 and 20 Julia threads. The tested
active environment used \texttt{DifferentialEquations} 7.16.1,
\texttt{PolyLog} 2.6.0, and \texttt{QuadGK} 2.11.2. The standalone directory
does not contain a project or manifest file, so the direct package versions
are specified explicitly. From the repository root:
\begin{verbatim}
cd "Numerical RG"
julia -e "using Pkg; for (name,version) in [
    ("DifferentialEquations","7.16.1"),
    ("PolyLog","2.6.0"), ("QuadGK","2.11.2")
]; Pkg.add(name=name,version=version); end"
julia --threads=auto test/runtests.jl
\end{verbatim}

The numerical values in the tables are reproduced with:
\begin{verbatim}
julia --threads=auto test/analytic_validation.jl
julia --threads=auto test/validation_metrics.jl
julia test/momentum_sum_check.jl
julia test/paper_moment_check.jl
\end{verbatim}

The grid-refinement and closure calculations are run with:
\begin{verbatim}
julia --threads=auto test/analytic_validation.jl 500 1000 2000 1223:300
\end{verbatim}

The asserted regression suite contains 542 checks. The repository workflow
\path{.github/workflows/numerical-rg.yml} installs the required packages and
runs the suite on Julia 1.11.

\begin{thebibliography}{99}
%\cite{Dixon:2019uzg}
\bibitem{Dixon:2019uzg}
L.~J.~Dixon, I.~Moult and H.~X.~Zhu,
%``Collinear limit of the energy-energy correlator,''
Phys. Rev. D \textbf{100} (2019) no.1, 014009
doi:10.1103/PhysRevD.100.014009
[arXiv:1905.01310 [hep-ph]].
%204 citations counted in INSPIRE as of 19 Jul 2026

\end{thebibliography}

\end{document}
